\chapter{Efficiency: Real Per-Sample Timing, GPU-Compute vs.\ Wall-Clock (Phase 2.6)}
\label{chap:p2-6}

\section{Goal}

The pilot's efficiency numbers were built from a 15-sample calibration and a linear extrapolation, with
no separation of GPU-compute time from end-to-end wall-clock and no standardized timing boundary. The
plan (item 4.6) asks for real per-sample timing with an explicit boundary (generate-only, device
synchronized, warm-up discarded), a GPU-compute-vs-wall-clock split to prove the workload is GPU-bound,
and --- deferred to the cluster --- a batched re-benchmark at multiple scales and a re-cost after the
gradient-attribution (D1) and dense-captioning (D4) decisions.

\section{What Was Done}

Two utilities were added. \texttt{summarize\_timings} discards a configurable number of warm-up
measurements (cold caches, kernel autotuning) and reports mean, median, and p90, so tail cost is visible
rather than hidden in a single calibrated mean. \texttt{time\_call\_cuda} wraps a call with a
standardized boundary: \texttt{torch.cuda.synchronize()} before and after, CUDA events to measure the
GPU-busy time, and \texttt{perf\_counter} for total wall time. The GPU time is the interval the GPU is
actually executing kernels; the wall time is everything. Their ratio is the GPU-bound-ness. Script~08
gained a \texttt{-{}-benchmark} mode that pre-fetches images (so I/O is outside the timed region ---
the local parquet cache makes this valid), discards warm-up samples, times each real sample, and reports
the split; the frozen calibration path is unchanged. Three unit tests cover the warm-up discarding and
p90.

\section{Results}

Timing 60 real samples (after 5 discarded warm-up samples), with the standardized boundary:

\begin{center}
\begin{tabular}{p{5.2cm}p{2.2cm}p{2.2cm}p{2.2cm}}
\toprule
\textbf{measure} & \textbf{mean} & \textbf{median} & \textbf{p90} \\
\midrule
GPU-compute time (ms/sample)    & 264.7 & 254.0 & 299.5 \\
end-to-end wall time (ms/sample) & 264.8 & 254.1 & 299.9 \\
\bottomrule
\end{tabular}
\captionof{table}{Real per-sample timing under the standardized GPU-synced boundary. GPU-compute and
wall-clock are essentially identical.}
\label{tab:p26-timing}
\end{center}

The GPU-compute-to-wall-clock ratio is \textbf{1.00}: the CUDA-event GPU time and the total wall time
coincide to within measurement noise. With images pre-fetched from the local cache and excluded from the
timed region, essentially all per-sample cost is GPU kernel execution --- there is no meaningful
CPU, Python, or I/O overhead left to shave. This is the direct, measured confirmation of the claim the
pilot could only assert: the pipeline is GPU-bound, and the local data cache (Phase~3) has removed the
streaming I/O that would otherwise contend. The real per-sample figure, 264.7~ms mean with a 299.5~ms
p90, replaces the pilot's calibrated point estimate and makes the tail explicit; at this cost one metric
pass over the full 3{,}004-image test split is about 13 GPU-minutes, extrapolated from a measured
per-sample number rather than a 15-sample calibration.

\section{What Is Deliberately Deferred}

Two of the plan's efficiency asks belong to hardware this pilot does not have and decisions not yet
taken, and are recorded as scale-run steps rather than forced here. The batched, multi-scale re-benchmark
(1k / 2.5k / 10k on the target cluster, to expose memory fragmentation and thermal effects that a
single-GPU per-sample number cannot) requires the cluster. And the re-cost after D1 (gradient attribution,
which the plan estimates at $10$--$50\times$ the forward-pass cost) and D4 (dense captioning, whose longer
outputs make the decode loop no longer free) cannot be measured until those decisions are made and their
code exists. The measurement machinery built here --- the synchronized boundary, warm-up discarding, and
the GPU-vs-wall split --- is exactly what those re-benchmarks will reuse; only the workload changes.

\section{Standard vs. Non-Standard Choices}

\textbf{Standard.} Synchronizing the device around a timed region, discarding warm-up, and reporting a
tail percentile are basic GPU-benchmarking hygiene that the pilot's calibration skipped.

\textbf{Non-standard, and the contribution.} Reporting GPU-compute and wall-clock as \emph{separate}
measured quantities, and using their ratio as an explicit GPU-bound-ness test, turns ``we believe it is
GPU-bound'' into a number (1.00). It also draws the line honestly at what a single-GPU pilot can measure:
the per-sample cost is real and reported, while the cluster-scale and post-D1/D4 costs are named as
pending rather than extrapolated from a machine and a workload that do not yet exist.

\section{Implications}

Every metric script can now be timed at its true per-sample cost with a consistent boundary, and the
GPU-bound result means the scale run's budget is governed by GPU throughput and batching, not I/O ---
which is what makes the deferred cluster batched benchmark the right next measurement. This closes
Phase~2: all six metrics now carry validity corrections (rule-aware, worker-loss-corrected, size-invariant,
dose-response, multi-hazard) and a real efficiency instrument, on a pipeline that still reproduces the
frozen pilot byte-for-byte under its default flags.
